\section{Kerangka Pemikiran}
\TemplateTodo{Jelaskan hubungan antara masalah, teori, metode, data, dan keluaran penelitian. Bagian ini dapat dibantu dengan gambar alur konseptual.}

\begin{figure}[H]
\centering
\includegraphics[width=0.75\linewidth]{example-figure.pdf}
\caption{Alur kerangka pemikiran penelitian}
\label{fig:framework}
\end{figure}

Gambar~\ref{fig:framework} menunjukkan alur kerangka pemikiran penelitian. Data ulasan dari Google Play Store akan melalui tahap \emph{preprocessing}, kemudian diekstraksi fiturnya menggunakan TF-IDF, dan diklasifikasikan menggunakan Na\"ive Bayes Classifier. Hasil klasifikasi dievaluasi menggunakan \emph{confusion matrix} untuk memperoleh nilai akurasi, precision, recall, dan F1-score. Contoh hasil pengujian metrik dapat dilihat pada Tabel~\ref{tab:csv-example} yang menampilkan perbandingan nilai metrik sebelum dan sesudah penerapan sistem.

\subsection{Contoh Tabel dari Data CSV}

Template dapat menampilkan tabel yang diisi dari berkas CSV. Simpan data di folder \texttt{references/} dan gunakan paket \texttt{csvsimple}.

\begin{table}[H]
\centering
\caption{Contoh tabel dari data CSV}
\label{tab:csv-example}
\csvautotabular{references/example.csv}
\end{table}

Tabel~\ref{tab:csv-example} menampilkan contoh data dari berkas \texttt{references/example.csv}. Pastikan jumlah kolom di CSV sesuai dengan tampilan yang diinginkan.

